\section{Length-Aware Conservation and a Nonrenewal Application}\label{sec:bounded60}
The full-graph recognition theorem quantifies over all source-to-sink paths. Its condition is sufficient, but need not be necessary, when only paths with at most $h$ edges are permitted. For example, edges $s\to t$, $s\to u$, $u\to t$ of capacities $1,1,1$ fail full-graph conservation; when $h=1$, only the first path is allowed and its capacity is unambiguous. The following extension makes the path quantifier explicit.

\begin{theorem}[Recognition under an edge allowance]\label{thm:bounded60}
Let $G$ be an acyclic graph with nonnegative rational edge capacities, a sink $t$, sources $S$, and an edge allowance $h$. After replacing $h$ by $\min\{h,|V|-1\}$, construct vertices $(v,\ell)$ for $0\leq\ell\leq h$, edges $(u,\ell)\to(v,\ell+1)$ of capacity $C_{uv}$, and zero-capacity edges from $(t,\ell)$ to one artificial sink. Prune vertices not on a path from some $(s,0)$ to the artificial sink. Then the following are equivalent: every original source has a common capacity across all its admissible paths of at most $h$ edges; the retained layered graph has a suffix-capacity potential.

Recognition requires $O(h(|V|+|E|))$ arithmetic operations and, upon failure, returns two original paths of at most $h$ edges from the same source with different capacities. Denote the source potentials by $H_{s,0}$. A complement resource with one source-independent terminal level exists exactly when $H_{s,0}-B_s=R$ is constant across admissible sources. Under the kernel assumptions of Theorem~\ref{thm:abstractapprox59}, the same feasible additive approximation holds with arithmetic count
$O(h|E|Q\log(Q+1)+h|V|Q)$, apart from kernel evaluation, where $Q=\lfloor2KhR/\varepsilon\rfloor+1$. The artificial edges are not counted or rounded.
\end{theorem}
\begin{proof}
A path to the artificial sink is in bijection with an original admissible path and its elapsed edge counts. The artificial edge adds no capacity. Apply Theorem~\ref{thm:recognition59} to the retained layered graph. A disagreement has two explicit layered suffixes and a source prefix. Erasing layer labels and the artificial edge gives the promised witness. Storing parent and successor pointers, and materializing paths only when a disagreement occurs, gives the linear bound in the lifted graph size; storing a separate full suffix at every vertex would not.

For any source, the capacity sum telescopes to $H_{s,0}$. Thus $w_e=C_e-x_e$ transforms its equality into $\sum w_e=H_{s,0}-B_s$. Constancy of this quantity is necessary and sufficient for the stated complement coordinate. Rounding an optimal admissible path down changes at most $h$ real edges. The same-path greedy repair is possible because its total capacity is at least $R$. The proof of Theorem~\ref{thm:abstractapprox59} therefore applies without change to the error bound. Each pair consisting of an original edge and elapsed edge count is processed once; the layer coordinate already used by the dynamic program is not duplicated. This proves the stated, rather than an extra factor-$h$, count.
\end{proof}

\paragraph{Oracle and bit model.} For general continuous kernels, the arithmetic counts are oracle counts: every requested rational grid value must be supplied exactly, and the theorem does not presume a polynomial-time implementation for an arbitrary continuous function. For rational piecewise-linear or piecewise-quadratic kernels supplied explicitly with rational breakpoints and polynomially encoded coefficients, exact evaluation at a grid point has polynomial bit cost in the input description and $\log(Q+1)$. The sum of at most $h$ such values has polynomially bounded encoding length. Multiplying the stated arithmetic count by this evaluation and integer-arithmetic cost gives a bound polynomial in the numerical $Q$ and the binary input length. An exponentially large grid is not made polynomial by encoding its size in binary. The renewal kernels have the separate constructive bound of Proposition~\ref{prop:cost56}.

\subsection{Corridor rehabilitation with endogenous contract boundaries}\label{sec:corridor60}
Consider an ordered set of admissible mileposts, with position $p(v)$, first post $s$, and final post $t$. An admissible edge $u\to v$ is a rehabilitation contract spanning the interval between the two posts. Its physical length is $C_{uv}=p(v)-p(u)$. A design chooses a compatible sequence of contracts partitioning the corridor, with at most $h$ contracts. On each chosen contract it treats a length $x_{uv}\in[0,C_{uv}]$, earns a continuous concave incremental condition benefit $\phi_{uv}(x_{uv})$, and pays a fixed mobilization charge included once in that edge payoff. A funding requirement fixes the total treated length at $B$.

This is precisely the resource-path model: vertices are permissible boundaries, edges are admissible contracts, the command-count analogue is the contract allowance, and $R=p(t)-p(s)-B$ is the untreated length. The potential is $H_v=p(t)-p(v)$, so all path capacities equal the corridor length regardless of the chosen partition. Deficit repair expands untreated lengths within the same contracts, keeping both the partition and the funding equality. Contract-specific technologies can change $\phi_{uv}$ without changing recognition. With several permissible starting posts, a common untreated length is possible exactly when $p(t)-p(s)-B_s$ is independent of $s$.

The mapping assumes divisible treatment and separable contract returns; it does not cover inter-contract deterioration spillovers or indivisible work packages. It is a decision model derived from stated physical quantities, not a field-calibrated experiment. It illustrates why conservation and recovery are useful outside renewal commands without importing the renewal-specific lottery or eligibility argument into an unrelated application.
